\documentclass[11pt,a4paper]{article}

\usepackage[utf8]{inputenc}
\usepackage[T1]{fontenc}
\usepackage[margin=2.2cm]{geometry}
\usepackage{parskip}
\usepackage{url}
\usepackage{enumitem}
\setlist{nosep,leftmargin=1.4em}

\pagestyle{plain}

\begin{document}

\begin{center}
{\large\textbf{Information Sheet}}\\[3pt]
``When Tool Calls Succeed but Workflows Fail: Anomalies at the
Agent--Tool Boundary''\\
Artem Trofimov and Boris Novikov
\end{center}

\subsection*{1. What is the main claim, and why is it an important
contribution to the autonomous agents and multiagent systems literature?}

LLM-based agents now act on the world through tools supplied by
independent parties: they send messages, place orders, and charge cards.
Our claim is that the correctness of such an agent cannot be settled by
the agent or its runtime alone. Whether its surviving external effects are
consistent with what the workflow decided depends on what the tools
declare at the boundary, and the standard annotations of the most widely
used tool interface cannot express what is needed.

The contribution to the agents literature is a specification-level
account of what it means for an agent's actions to be correct when they
are retried, executed speculatively, run concurrently with other agents,
and observed by outside parties. We give an effect-history model that
separates events in the world from the agent's observations of them; a
catalog of eight anomalies, each a forbidden pattern in that history,
grouped by outcome uncertainty, workflow structure, and interaction with
other agents or exogenous observers; the boundary capability needed to
exclude each anomaly; and four limits beyond which no runtime can provide
a guarantee
without additional boundary support, such as authoritative tool
participation, shared mediation, or visibility control. In multiagent-systems terms,
a tool is an independently operated service outside the agent's control,
whose call interface need not define a shared protocol for resolving
outcomes or managing external effects. Work on information protocols and
commitments shows how fault tolerance follows once all parties enact a
common protocol~[1,2,3]. Our capabilities state what a tool would have to
declare to become such a participant.

\subsection*{2. What evidence supports the claim?}

\begin{itemize}
\item \emph{Analytical.} The model, the catalog, and the capability
derivation are the main results. The catalog is organized by
construction: the uncertainty anomalies follow from the three actions a
runtime can take under an unresolved outcome (re-issue, commit,
compensate), the workflow anomalies from the relations of an effect to its
resolution point, and the interaction anomalies from the two kinds of
party an effect can meet. The four guarantee limits are argued from the
definitions rather than proved: two correspond to classical results, the
exactly-once barrier and atomic commitment, and formal proofs for all
four are left as future work.
\item \emph{Comparative.} We place six recent agent runtimes, including
Atomix~[4], Cordon~[5], and CoAgent~[6], against the catalog. Each covers a
fragment, and each obtains the tool properties it needs out of band,
through its own adapter declarations.
\item \emph{Empirical.} A census of 98,291 tools from the official Model
Context Protocol registry shows that annotation fields are widely emitted
but carry only coarse call-level hints. None of the eight capabilities is
fully expressible in the standard vocabulary. The census supports this
claim about what the interface declares; it does not, and is not meant
to, establish the guarantee limits. Data, code, and methodology are
public.
\end{itemize}

\subsection*{3. Which papers by other authors make the most closely related
contributions, and how is this paper related to them?}

\begin{itemize}
\item \emph{Fault-tolerant interaction protocols}~[1,2,3]. Mandrake
deduplicates repeated messages through key-based correlation when every
party enacts the same information protocol, a protocol-level analogue of
the logical-operation identity that A1 requires. We study the boundary where the
other party is an independently supplied tool that cannot be assumed to
enact such a protocol, and derive what it would need to expose.
\item \emph{Transactional agent runtimes}~[4,5,6,7]. These systems
implement staging, compensation, or concurrency control for tool calls. We
do not propose another runtime; we characterize which guarantees any
runtime can establish from what the boundary declares, and read these
systems as partial coverage of the catalog.
\item \emph{Anomaly catalogs for agent systems}~[8]. A recent verified
catalog of concurrency anomalies targets a runtime's shared state. Ours
centers the outcome uncertainty of calls to third-party tools and the
interface capabilities that exclude each anomaly.
\item \emph{Transaction processing}~[9,10,11]. We borrow anomaly-based
correctness and multilevel derivation of guarantees, and show that the
assumption those models make, that operation semantics are known, fails at
the agent--tool boundary.
\end{itemize}

\subsection*{4. Have parts of this paper been published elsewhere?}

No part of this paper has appeared in an archival venue. A preprint is
available at \url{https://arxiv.org/abs/2609.15397}, and an earlier
version has been accepted at ``Who Verifies the Agents?'', a non-archival
NeurIPS 2026 workshop without proceedings, to be held in December 2026.

\subsection*{References}

\small
\begin{enumerate}[label={[\arabic*]},leftmargin=2em]
\item S.~H. Christie V, A.~K. Chopra, M.~P. Singh. Mandrake: multiagent
systems as a basis for programming fault-tolerant decentralized
applications. \emph{Autonomous Agents and Multi-Agent Systems}
36(1):16, 2022.
\item M.~P. Singh. Information-driven interaction-oriented programming:
BSPL, the blindingly simple protocol language. In \emph{Proc.\ AAMAS},
pp.~491--498, 2011.
\item P. Yolum, M.~P. Singh. Commitment machines. In \emph{Intelligent
Agents VIII (ATAL 2001)}, LNAI 2333, pp.~235--247, Springer, 2002.
\item B. Mohammadi et al. Atomix: Timely, transactional tool use for
reliable agentic workflows. arXiv:2602.14849, 2026.
\item Z. Chen et al. Cordon: Semantic transactions for tool-using LLM
agents. arXiv:2606.17573, 2026.
\item H. Lyu et al. CoAgent: Concurrency control for multi-agent systems.
arXiv:2606.15376, 2026.
\item S. Perera, K. Hapuarachchi, F. Leymann, R. Khalaf. Robust Agent
Compensation (RAC): Teaching AI agents to compensate. In \emph{Proc.\ ACM
Conference on AI and Agentic Systems}, pp.~253--262, 2026.
\item S. Khan. Verified detection and prevention of concurrency anomalies
in multi-agent large language model systems. arXiv:2606.17182, 2026.
\item G. Weikum. Principles and realization strategies of multilevel
transaction management. \emph{ACM Transactions on Database Systems}
16(1):132--180, 1991.
\item H. Garcia-Molina, K. Salem. Sagas. In \emph{Proc.\ ACM SIGMOD},
pp.~249--259, 1987.
\item H. Berenson et al. A critique of ANSI SQL isolation levels. In
\emph{Proc.\ ACM SIGMOD}, pp.~1--10, 1995.
\end{enumerate}

\end{document}